\section{Proposed Approach} % present algorithm
\label{sec:proposed_approach}

\subsection{Complete Architectural Diagram}
\label{subsec:architecture}

Inside each node, the health-checking process follows an event-driven structure (1) \& (2.b). A node regularly checks
its neighbors and reacts when abnormal behavior is detected, such as missing responses, delayed replies, or unreachable
nodes. When such an event occurs, the node may request indirect checks from other neighbors, calculate the suspicion
level of the target node, and categorize it. After the status is determined, the result will be broadcasted to the
system so that the cluster can maintain a consistent view of node conditions.

The lower part of the diagram (3) represents the supporting observation subsystem. This subsystem receives node status
information from the decentralized cluster and passes it through several logical layers. The first layer (3.a) receives
reported node status, the next layer (3.c) processes and organizes the information, and the storage layer (3.c) \&
(3.d) keeps historical records of node states and failure events. The visualization layer (3.b) then retrieves this
information and presents it in a readable form for users to inspect the current and past condition of the cluster.

The decentralized cluster (1) is responsible for detection, classification, and reaction, while the supporting layers
(3) are used only for reporting, storage, and visualization. Therefore, the system can continue checking node health
and maintaining neighbor status even if the observation subsystem, storage component, or visual interface becomes
unavailable.

Fig.~\ref{fig:architecture} shows the complete architecture of the proposed system.

% \begin{figure*}[t]
%     \centering
%     
\includegraphics[width=0.9\textwidth]{Sample.jpeg}
%     \caption{Complete architectural diagram of the node health-checking system.}
%     \label{fig:architecture}
% \end{figure*}

\begin{figure*}[!t]
  \centering
  \includegraphics[width=0.95\textwidth]{ArchitecturalDiagram.png}
  \caption{Complete architectural diagram of the node health-checking system.}
  \label{fig:architecture}
\end{figure*}

% \resizebox{0.95\textwidth}{!}{%
% \begin{circuitikz}
% \tikzstyle{every node}=[font=\fontsize{18.2pt}{23.7pt}\selectfont]
% \draw [fill=white, fill opacity=1, line width=0.2pt] (-2.5,17.25) rectangle (38.75,-6.5);

% \draw  (1.25,13.5) circle (1.125cm);
% \draw  (5,7.25) circle (1.125cm);
% \draw  (6.25,14.75) circle (1.125cm);
% \draw  (1.25,9.75) circle (1.125cm);
% \node [font=\fontsize{19.9pt}{25.9pt}\selectfont, inner xsep=0.080cm, inner ysep=0.085cm, rounded
% corners=0.020cm] at (1.25,13.5) {node};
% \node [font=\fontsize{19.9pt}{25.9pt}\selectfont, inner xsep=0.080cm, inner ysep=0.085cm, rounded
% corners=0.020cm] at (1.25,9.75) {node};
% \node [font=\fontsize{19.9pt}{25.9pt}\selectfont, inner xsep=0.080cm, inner ysep=0.085cm, rounded
% corners=0.020cm] at (5,7.25) {node};
% \node [font=\fontsize{19.9pt}{25.9pt}\selectfont, inner xsep=0.080cm, inner ysep=0.085cm, rounded
% corners=0.020cm] at (6.25,14.75) {node};
% \draw [short] (2.375,13.5) -- (6.25,13.625);
% \draw [short] (6.25,13.625) -- (5,8.375);
% \draw [short] (2.375,13.5) -- (2.375,9.75);
% \draw [short] (2.375,9.75) -- (5,8.375);
% \draw [-{Stealth[scale=1.5]}, ] (6.25,13.625) -- (10,8.5);
% \draw [ rounded corners = 37.5] (10,9.75) rectangle (15,7.25);
% \draw [-{Stealth[scale=1.5]}, ] (2.375,13.5) -- (10,8.5);
% \draw [-{Stealth[scale=1.5]}, ] (5,8.375) -- (10,8.5);
% \draw [-{Stealth[scale=1.5]}, ] (2.375,9.75) -- (10,8.5);
% \node [font=\fontsize{18.2pt}{23.7pt}\selectfont, fill={rgb,255:red,255; green,255; blue,255}, fill
% opacity=1, text opacity=1, inner xsep=0.080cm, inner ysep=0.085cm, rounded corners=0.020cm] at (12.5,8.5) {SERVER};
% \draw [ rounded corners = 37.5] (17.5,2.25) rectangle (22.5,-0.25);
% \draw [ rounded corners = 37.5] (25,2.25) rectangle (30,-0.25);
% \draw [ rounded corners = 37.5] (32.5,2.25) rectangle (37.5,-0.25);
% \draw [line width=0.2pt, -{Stealth[scale=1.5]}, ] (12.5,7.25) -- (17.5,1);
% \draw [line width=0.2pt, -{Stealth[scale=1.5]}, ] (22.375,1.625) -- (25.125,1.625);
% \draw [line width=0.2pt, -{Stealth[scale=1.5]}, ] (29.875,1.625) -- (32.625,1.625);
% \draw [line width=0.2pt, -{Stealth[scale=1.5]}, ] (32.625,0.375) -- (29.875,0.375);
% \draw [line width=0.2pt, -{Stealth[scale=1.5]}, ] (25.125,0.375) -- (22.375,0.375);
% \node [font=\fontsize{19.9pt}{25.9pt}\selectfont, inner xsep=0.080cm, inner ysep=0.085cm, rounded
% corners=0.020cm] at (19.875,1) {Controller};
% \node [font=\fontsize{19.9pt}{25.9pt}\selectfont, inner xsep=0.080cm, inner ysep=0.085cm, rounded
% corners=0.020cm] at (27.375,1) {Service};
% \node [font=\fontsize{19.9pt}{25.9pt}\selectfont, inner xsep=0.080cm, inner ysep=0.085cm, rounded
% corners=0.020cm] at (35,1) {Repository};
% \draw [ rounded corners = 37.5] (17.5,14.75) rectangle (22.5,12.25);
% \node[anchor=center, align=center, fill={rgb,255:red,255; green,255; blue,255}, fill opacity=1, text
% opacity=1, inner xsep=0.080cm, inner ysep=0.085cm, rounded corners=0.020cm] at (20,13.5)
% {\fontsize{18.2pt}{23.7pt}\selectfont Dashboard};
% \draw [line width=0.2pt, -{Stealth[scale=1.5]}, ] (20.625,12.25) -- (20.625,2.25);
% \draw [line width=0.2pt, -{Stealth[scale=1.5]}, ] (19.375,2.25) -- (19.375,12.25);
% \draw [ fill={rgb,255:red,255; green,255; blue,255}, fill opacity=1, line width=0.2pt , rounded corners =
% 37.5, dashed] (17.5,-2.75) rectangle (22.5,-5.25);
% \node[anchor=center, align=center, fill={rgb,255:red,255; green,255; blue,255}, fill opacity=1, text
% opacity=1, inner xsep=0.080cm, inner ysep=0.085cm, rounded corners=0.020cm] at (20,-4)
% {\fontsize{18.2pt}{23.7pt}\selectfont collect};
% \draw [ fill={rgb,255:red,255; green,255; blue,255}, fill opacity=1, line width=0.2pt , rounded corners =
% 37.5, dashed] (25,-2.75) rectangle (30,-5.25);
% \node[anchor=center, align=center, fill={rgb,255:red,255; green,255; blue,255}, fill opacity=1, text
% opacity=1, inner xsep=0.080cm, inner ysep=0.085cm, rounded corners=0.020cm] at (27.5,-4)
% {\fontsize{18.2pt}{23.7pt}\selectfont process};
% \draw [ rounded corners = 37.5] (32.5,-2.75) rectangle (37.5,-5.25);
% \node [font=\fontsize{19.9pt}{25.9pt}\selectfont, inner xsep=0.080cm, inner ysep=0.085cm, rounded
% corners=0.020cm] at (35,-4) {Database};
% \draw [line width=0.2pt, -{Stealth[scale=1.5]}, dashed] (20,-2.75) -- (20,-0.25);
% \draw [line width=0.2pt, -{Stealth[scale=1.5]}, dashed] (27.5,-2.75) -- (27.5,-0.25);
% \draw [line width=0.2pt, -{Stealth[scale=1.5]}, ] (35.625,-0.25) -- (35.625,-2.75);
% \draw [line width=0.2pt, -{Stealth[scale=1.5]}, ] (34.375,-2.75) -- (34.375,-0.25);
% \node [font=\fontsize{18.2pt}{23.7pt}\selectfont, fill={rgb,255:red,255; green,255; blue,255}, fill
% opacity=1, text opacity=1, inner xsep=0.080cm, inner ysep=0.085cm, rounded corners=0.020cm, align=center]
% at (11.125,3.75) {node-info.json \\ ~ \\ POST METHOD};
% \node [font=\fontsize{18.2pt}{23.7pt}\selectfont, fill={rgb,255:red,255; green,255; blue,255}, fill
% opacity=1, text opacity=1, inner xsep=0.080cm, inner ysep=0.085cm, rounded corners=0.020cm, align=center]
% at (24.125,8.375) {request node-info.json \\ ~ \\ GET METHOD};
% \end{circuitikz}
% }%

% NOTE: add something to lead to these algo, algo will be divided into three subsections
\subsection{Failure Detection: Hybrid SWIMPHI}
\label{subsec:hybrid_swimphi}
The proposed node health-checking system is influenced by both SWIM and the $\phi$ accrual failure detector.
From SWIM, the system adopts the idea of decentralized peer-to-peer checking and helper-based indirect probing.
This allows more than one node to participate in checking the status of a target node and reduces dependence
on a single monitoring component.
From the $\phi$ accrual detector, the system adopts the idea of using timing history and a continuous
suspicion value instead of making an immediate binary decision after one missed response.
In the proposed design, a target node is first checked through heartbeat or ping messages.
If the target does not respond, helper nodes from the neighbor list can be used to verify whether the problem
is a real node failure or only a local communication problem.
After that, the $\phi$ value is used to evaluate the seriousness of the failure suspicion based on the
response-time history.

By combining these two ideas, the proposed approach aims to improve both scalability and accuracy. SWIM contributes a
scalable architecture for distributed membership checking, while the $\phi$ accrual detector contributes an adaptive
decision mechanism for handling delay and message loss. Therefore, the system can avoid immediately deleting a node
after a single missed heartbeat and can instead move through health states such as healthy, suspected, and failed based
on both cooperative probing and suspicion-level evaluation.

As illustrated in Fig.~\ref{fig:phi_pipeline}, the algorithm combines indirect heartbeat verification with the $\phi$
accrual failure detector. The indirect heartbeat mechanism reduces false-positive failure detection by allowing
neighbor nodes to verify a suspicious node, while the $\phi$ accrual detector provides an adaptive suspicion value
based on recent response-time behavior.

\begin{figure*}[!t]
  \centering
  \includegraphics[width=0.95\textwidth]{FailureDetectionProcedure.png}
  \caption{Pipeline of indirect heartbeat verification and $\phi$ accrual failure detection.}
  \label{fig:phi_pipeline}
\end{figure*}

\subsubsection{Indirect Heartbeat Verification}
\label{subsubsec:indirect_heartbeat_verification}

Each node maintains a neighbor list containing the nodes with which it has a direct bidirectional connection.

When a node wants to check the availability of one of its neighbors, it sends a ping message and waits for an
acknowledgement. For example, if node $x_2$ wants to check node $x_1$, node $x_2$ sends a direct ping request to node
$x_1$. If node $x_1$ replies within the expected time range, node $x_2$ records the response time and classifies node
$x_1$ as reachable. This successful response also becomes part of the historical response-time data used later by the
$\phi$-accrual detector.

However, if node $x_2$ does not receive an acknowledgement, it does not immediately mark node $x_1$ as failed. Instead,
the system uses indirect heartbeat verification. Specifically, node $x_2$ requests other nodes in its neighbor list to
verify node $x_1$, such as nodes $x_3$ and $x_5$, to ping node $x_1$ on its behalf. If any helper node receives an
acknowledgement from node $x_1$, it notifies node $x_2$, and node $x_1$ is treated as alive. In this case, the problem
is considered a possible connection issue between nodes $x_1$ and $x_2$ rather than a complete node failure.

If neither the direct ping nor the indirect pings receive a response from node $x_1$, the system has stronger evidence
that node $x_1$ may be unreachable. At this stage, the system does not rely only on the missing acknowledgements.
Instead, it applies the $\phi$-accrual failure detector to calculate a suspicion value based on recent response-time
history.

\subsubsection{\texorpdfstring{$\phi$}{phi} Accrual Calculation}
\label{subsubsec:phi_accrual_calculation}

The $\phi$ accrual detector avoids making a binary alive-or-dead decision based solely on a fixed timeout. Instead, it
calculates a continuous suspicion value, $\phi$.

To make this decision, the detector relies on historical response-time data to determine whether the current waiting
time is normal or abnormal. Each time a ping is sent and an acknowledgement is successfully received, the system
records the response time. These recorded samples represent the recent communication behavior between nodes.

To keep the detector adaptive, response-time samples are stored in a sliding window. This window operates like a queue:
when a new response-time sample is recorded, it is added to the window. If the window has already reached its maximum
size, the oldest sample is removed. As a result, the window always contains the most recent response-time samples.

Let the response-time samples in the sliding window be:

\begin{equation}
  X = \{x_1, x_2, x_3, \ldots, x_n\}
\end{equation}

The mean response time and standard deviation are calculated as:

\begin{equation}
  \mu = \frac{1}{n}\sum_{i=1}^{n}x_i
\end{equation}

\begin{equation}
  \sigma = \sqrt{\frac{1}{n}\sum_{i=1}^{n}(x_i - \mu)^2}
\end{equation}

where $n$ is the number of samples, $x_i$ is a recorded response time, $\mu$ is the average response time, and $\sigma$
represents response-time variation. These values are then used to estimate the response-time distribution for the
$\phi$ calculation.

After the response-time distribution is estimated from the sliding window, the system calculates how unusual the
current waiting time is. Let $t$ be the current waiting time since the ping was sent. The probability that a normal
response time is greater than $t$ is defined as:

\begin{equation}
  P(T > t) = 1 - F(t)
\end{equation}

where $F(t)$ is the cumulative distribution function of the response-time distribution. The suspicion value is then
calculated as:

\begin{equation}
  \phi(t) = -\log_{10}(P(T > t)) = -\log_{10}(1 - F(t))
\end{equation}

A higher $\phi$ value means that the current waiting time is more unusual compared with recent response-time history,
so the target node is more likely to be unreachable.

\subsubsection{Node State Classification}
\label{subsubsec:node_state_classification}

The final output of the methodology is the classification of the monitored node. After the $\phi$ value is calculated,
the system compares it with predefined thresholds. These thresholds are configurable by the developer depending on how
sensitive the failure detection process should be. A lower threshold can detect unreachable nodes faster, but may
increase false positives. A higher threshold can reduce false positives, but may delay failure classification.

\begin{table}[h]
  \centering
  \caption{Node state classification based on $\phi$ value}
  \label{tab:phi_classification}
  \begin{tabular}{|c|c|}
    \hline

    \hline

    \hline

    \hline

    \hline

    \hline
  \end{tabular}
\end{table}

Through this combination of indirect verification and adaptive suspicion calculation, the proposed method can reduce
false-positive failure detection caused by temporary delay, message loss, or unstable direct connections.

\subsection{Neighbor Discovery: FINDING NEMO}
\label{subsec:node_discovery}

FINDING NEMO, short for \textit{Finding Neighbors and Evicting Members of their Own}, is the neighbor discovery and
network joining algorithm proposed in this work. The name reflects the fact that a joining node must identify both
prospective neighbors and any existing members that are evicted by those neighbors during the joining process, as
discussed later in this section. To facilitate understanding of the algorithm, several terms and system rules are first
defined as follows:

\begin{enumerate}
  \item \textbf{Neighbor}: A bidirectional relationship established between two nodes in the network. If
        node A is a neighbor of node B, then node B is also considered a neighbor of node A.

  \item \textbf{Neighbor List}: A local data structure maintained by each node that stores information
        about its neighbors, such as their network addresses. The neighbor list has a fixed capacity of $k$,
        where $k$ is an even number. To maintain network connectivity and resilience, each node attempts to
        keep its neighbor list as full as possible whenever sufficient peers are available.

\end{enumerate}

FINDING NEMO begins when a newly created node requests to join the network. To promote random neighbor selection, the
joining node broadcasts a join request to all existing nodes. Assuming ideal network connectivity, two cases are
considered:

\begin{enumerate}

  \item \textbf{Small Network.}
        If the network contains at most $k+1$ nodes, the joining node establishes connections with all existing nodes.
        This allows the node to reach the desired neighbor count without requiring any restructuring of the
        network topology.

        Figure~\ref{fig:small_network} illustrates this scenario. Consider a network consisting of four nodes and a target
        neighbor count of $k=4$. When node $x_5$ joins the network, it establishes connections with all existing nodes, thereby
        satisfying its neighbor requirement.
        \begin{figure}[H]
          \centering
          \includegraphics[width=0.3\textwidth]{FINDING_NEMO_small_pretty.png}
          \caption{FINDING NEMO in a small network}
          \label{fig:small_network}
        \end{figure}

  \item \textbf{Scaled Network.}
        If the network contains more than $k+1$ nodes, all existing nodes are assumed to have already reached
        their target neighbor count.
        Nevertheless, they still acknowledge the join request.
        After collecting all acknowledgements, the joining node constructs a join plan by randomly selecting
        $k/2$ responding nodes.
        The joining node immediately establishes connections with these selected nodes.
        Subsequently, each selected node is instructed to evict one of its current neighbors.
        The evicted nodes then establish connections with the joining node.
        Consequently, the joining node acquires $k/2$ neighbors from the selected nodes and an additional
        $k/2$ neighbors from the evicted nodes, resulting in a total of $k$ neighbors.

        Figure~\ref{fig:scaled_network} presents an example of this process. Building upon the network shown in
        Figure~\ref{fig:small_network}, assume that node $x_6$ requests to join a network of five nodes with $k=4$. One
        possible outcome of the random selection process is that nodes $x_5$ and $x_3$ are selected. Node $x_5$ subsequently
        evicts node $x_2$, while node $x_3$ evicts node $x_4$. Node $x_6$ therefore establishes direct connections with $x_5$
        and $x_3$, as well as with the evicted nodes $x_2$ and $x_4$. As a result, node $x_6$ obtains four neighbors while
        preserving the desired network density.
        \begin{figure}[H]
          \centering
          \includegraphics[width=0.3\textwidth]{FINDING_NEMO_scaled_pretty.png}
          \caption{FINDING NEMO in a scaled network}
          \label{fig:scaled_network}
        \end{figure}

\end{enumerate}

This process maintains a bounded network density while preserving randomness in neighbor assignment, thereby reducing
the likelihood of node isolation and structural bias within the topology. In practice, a hybrid selection strategy is
employed. Nodes whose neighbor lists are not yet full are prioritized during the acknowledgement phase. By
preferentially connecting to such nodes, the algorithm accelerates the completion of underpopulated neighbor lists
while remaining consistent with the network's connectivity constraints.

\subsection{Connectivity Healing through Under-connected and Adaptive Link Adjustment for Network Health
  (CHUALANH) algorithm:}
\label{subsec:node_deletion}

\subsubsection{\textbf{Topology Repair Mechanism}}
\label{subsubsec:repair_mechanism}
In a decentralized architecture, node failures may disrupt the desired network structure. Specifically, when
a node unexpectedly fails, all of its incident edges are deleted. Consequently, some neighboring nodes lose
one connection and their degree decreases from $k$ to $k - 1$. These nodes are referred to as \textbf{deficient nodes}

Since the system is designed to maintain $k$ neighbors for every node, a topology repair process is required to restore
the desired degree constraint. To address this problem, we propose a \textbf{Connectivity Healing through
  Under-connected and Adaptive Link Adjustment for Network Health (CHUALANH)} algorithm.

\subsubsection{\textbf{Deficient Nodes}}
\label{subsubsec:deficient_node}
Let $G = (V, E)$ denote the network topology, where each node maintains exactly $k$ neighbors. Assume that
node $v_f$ fails and is removed from the network.

All nodes previously connected to $v_f$ become deficient nodes. Let $A = \{a_1, a_2, \ldots, a_m\}$ denote the set of
deficient nodes. Each deficient node requires one additional connection to restore its degree to $k$

\subsubsection{\textbf{Repair Procedure}}
\label{subsubsec:repair_procedure}
The repair process repeatedly selects two deficient nodes, denoted by $x_1$ and $x_2$, and restores their
missing connection according to the following cases.

\begin{enumerate}
  \item \textbf{Case 1: No Existing Connection}

        If $x_1$ and $x_2$ are not directly connected, a new edge is eslablished between them: $$(x_1, x_2) \in E$$ As a
        result, both nodes recover one connection.

  \item \textbf{Case 2: Existing Connection}

        If an edge already exists between $x_1$ and $x_2$, a direct connection between is impossible because it would create a
        duplicate connection

        In this case, a connected pair of nodes $x_3$ and $x_4$ is randomly selected such that $$(x_3, x_4) \in E$$ The
        connection between $x_3$ and $x_4$ is removed and replaced with two new edges. Depending on the local topology, one of
        the following rewiring patterns is applied: $$(x_3, x_4) \rightarrow (x_1, x_3), (x_2, x_4)$$ or $$(x_3, x_4)
          \rightarrow (x_1, x_4), (x_2, x_3)$$ The selected pattern must ensure that no duplicate edges are created and that the
        degree of every involved node remains remains unchanged except for the deficient node $x_1$ and $x_2$, whose degrees
        increase by one.

\end{enumerate}

\subsubsection{\textbf{Topological Constraints}}
\label{subsubsec:topological_constraints}
Not every candidate pair $(x_3, x_4)$ is suitable for rewiring. Additional constraints are imposed to prevent
duplicate connections and preserve network consistency.

Consider four nodes $x_1$, $x_2$. $x_3$, $x_4$.

\begin{enumerate}
  \item \textbf{Scenario A: General Configuration}

        The four nodes may already contain several existing edges, including one or two diagonal connections. The rewiring
        operation is valid provided that the newly created edges do not duplicate any existing connections. In this case, the
        edge $(x_3, x_4)$ is removed and replaced with $$(x_1, x_3), (x_2, x_4)$$
        \begin{figure}[H]
          \centering
          \includegraphics[width=0.3\textwidth]{Rewiring1_pretty.png}
          \caption{Scenario A: topological constraint}

          \vspace{0.5cm}

          \includegraphics[width=0.3\textwidth]{Rewiring2_pretty.png}
          \caption{Scenario A after rewiring}
        \end{figure}

  \item \textbf{Scenario B: Rectangular Configuration}

        When the four nodes form a rectangular structure or contain only one existing vertical edges, the alternative rewiring
        pattern is applied: $$(x_3, x_4) \rightarrow (x_1, x_4), (x_2, x_3)$$ This operation restores the degrees of the
        deficient nodes while maintaining the degree of all other nodes.
        \begin{figure}[H]
          \centering
          \includegraphics[width=0.3\textwidth]{Rewiring3_pretty.png}
          \caption{Scenario B: topological constraint}

          \vspace{0.5cm}

          \includegraphics[width=0.3\textwidth]{Rewiring4_pretty.png}
          \caption{Scenario B after rewiring}
        \end{figure}
\end{enumerate}

\subsubsection{\textbf{Termination}}
\label{subsubsec: termination}
The repair process is repeated until every deficient nodes have recovered their missing connections and once
again possess $k$ neighbors. The parameter k is assumed to be a globally predefined even integer, ensuring
that the rewiring mechanism terminates when a neighbor list size of each deficient node reaches the value k.
\subsection{Handling False Positives}
\label{subsec:false_positives}
Although the proposed approach significantly reduces false-positive detections, rare cases may still occur
under severe network congestion or temporary network partitions, where the health-checking process is disrupted.
In such situations, a healthy node may be incorrectly classified as \textbf{UNREACHABLE}.

To maintain consistency of the network topology, nodes classified as \textbf{UNREACHABLE} are removed from the active
membership view and treated as departed members. Once connectivity is restored, the deficient node is required to
rejoin the network through the standard membership protocol, thereby rebuilding its neighbor relationships and
restoring its participation in the system.

In this prototype, automatic process recovery is considered a deployment-level concern rather than a responsibility of
the distributed protocol itself. The protocol only determines when a node should be considered unreachable and
subsequently rejoin the network. In a production environment, this mechanism can be complemented by external
supervisors such as systemd, Docker restart policies, Kubernetes controllers, or watchdog services to automatically
restart failed node processes.
